\subsection{Universally measurable functions}

In this issue, X is a locally compact topological space. We recall (INT, V, p. 28, § 3, n° 4, def. 2) that a map f from X into a topological space Y is universally measurable if it is \(\mu\) -measurable for every positive measure \(\mu\) on X. It suffices for this that f be \(\mu\) -measurable for every positive measure \(\mu\) with compact support on X (INT, V, p. 28, § 4, no. 3, prop. 6).

Lemma 4. --- Let Y and Z be locally compact topological spaces, \(f\colon X\to Y\) and \(g\colon Y\to Z\) universally measurable mappings. The mapping \(g\circ f\colon X\to Z\) is universally measurable.

Let \(\mu\) be a positive measure with compact support on X, and let C be its support. The restriction of \(f\) to C is \((\mu|\mathrm{C})\)-proper. The map \(g\) is measurable with respect to the image measure \((f|\mathrm{C})(\mu|\mathrm{C})\), so the map \((g \circ f)|\mathrm{C} = g \circ (f|\mathrm{C})\) is \((\mu|\mathrm{C})\)-measurable. Consequently, the map \(g \circ f\) is \(\mu\)-measurable. The lemma follows.

We denote by \(\mathcal{L}_{\mathrm{u}}(\mathrm{X})\) the vector space of complex-valued functions that are universally measurable on X and by \(\mathcal{L}_{\mathrm{u}}^{\infty}(\mathrm{X})\) the subspace of functions \(f \in \mathcal{L}_{\mathrm{u}}(\mathrm{X})\) which are bounded on X. For \(f \in \mathcal{L}_{\mathrm{u}}^{\infty}(\mathrm{X})\), we denote by \(\|f\|_{\infty} = \sup_{x \in X} |f(x)|\) .

Proposition 3. --- a) The set \(\mathcal{L}_{\mathrm{u}}(\mathrm{X})\) is an involutive subalgebra of the involutive algebra of functions from X to C;

\begin{enumerate}
\def\labelenumi{\alph{enumi})}
\setcounter{enumi}{1}
\item
  The set \(\mathcal{L}_{\mathrm{u}}^{\infty}(\mathrm{X})\) is a subalgebra of \(\mathcal{L}_{\mathrm{u}}(\mathrm{X})\) and the mapping defined by \(f\mapsto \| f\|_{\infty}\) is a norm on \(\mathcal{L}_{\mathrm{u}}^{\infty}(\mathrm{X})\) for which \(\mathcal{L}_{\mathrm{u}}^{\infty}(\mathrm{X})\) is a unital star algebra;
\item
  The algebra \(\mathcal{L}_{\mathrm{u}}(\mathrm{X})\) contains \(\mathcal{C}(\mathrm{X})\) and \(\mathcal{L}_{\mathrm{u}}^{\infty}(\mathrm{X})\) contains \(\mathcal{C}_b(\mathrm{X})\) ;
\item
  Every Borel function from X to \(\mathbf{C}\) (TG, IX, p. 61, def. 5) belongs to \(\mathcal{L}_{\mathrm{u}}(\mathrm{X})\) ;
\item
  For any sequence \((f_n)_{n\in \mathbf{N}}\) in \(\mathcal{L}_{\mathrm{u}}(\mathrm{X})\) which converges pointwise to a function \(f\) from X into \(\mathbf{C}\), we have \(f\in \mathcal{L}_{\mathrm{u}}(\mathrm{X})\) .
\end{enumerate}

The assertions \(a\) ) and \(c\) ) follow from the definitions (cf. INT, IV, p. 175, § 5, no. 3, cor. 3). Assertion \(d\) ) is a consequence of INT, IV, p. 179, § 5, no. 5, th. 4, since the inverse image under a Borel function of every Borel subset of \(\mathbf{C}\) is a Borel subset of X, which is universally measurable (INT, V, p. 28, § 3, no. 4). Finally, assertion \(e\) ) follows from Egoroff's theorem (INT, IV, p. 175, § 5, no. 4, th. 2).

To prove assertion \(b\) ), it suffices to note that assertion \(e\) ) implies, \(a\) fortiori, that a uniform limit of universally measurable functions is universally measurable.

There may exist universally measurable functions on X that are not Borel. Indeed, if X is metrizable, the characteristic function of a Souslin subset of X is universally measurable (cf. INT, IV, p. 171, § 5, no. 1, cor. 2); however, in every uncountable Polish space, there exists a Souslin subset that is not Borel ("Souslin's theorem"; this follows from TG, IX, p. 120, exerc. 8 and from the fact that for every uncountable Polish space X, there exists a Borel bijection \(\mathrm{X}\to\mathbf{N}^{\mathbf{N}}\) whose inverse is Borel).

Lemma 5. --- Let \(\mu\) be a positive measure on X. Let \(g\) be a \(\mu\)-measurable function from X into \(\mathbf{C}\) and let \(S\) be its \(\mu\)-essential image. For every function \(f \in \mathcal{L}_{\mathrm{u}}(\mathrm{S})\), the function \(h\) from X into \(\mathbf{C}\) such that \(h(x) = 0\) if \(g(x) \notin \mathrm{S}\) and \(h(x) = f(g(x))\) if \(g(x) \in \mathrm{S}\) is \(\mu\)-measurable.

Let \(x \in X\) and \(U\) be an open relatively compact neighborhood of x. The set \(\mathrm{N} = \mathrm{U} - (\mathrm{U} \cap g^{-1}(\mathrm{S}))\) is \((\mu|\mathrm{U})\)-negligible. For every integer \(n \geqslant 1\), let \(V_{n}\) be an open set such that \(N \subset V_{n} \subset U\) and \(\mu(V_{n} - N) < 1/n\) (INT, IV, p. 116, § 1, no. 4, prop. 19).

The restriction \(\widetilde{g}\) of \(g\) to \(\mathrm{U}-\mathrm{V}_n\) is \(\mu|(\mathrm{U}-\mathrm{V}_n)\)-proper and its image is contained in S. Since \(f\) is universally measurable, the map \(x\mapsto f(g(x))\) from \(\mathrm{U}-\mathrm{V}_n\) into \(\mathbf{C}\) is measurable with respect to the measure \(\mu|(\mathrm{U}-\mathrm{V}_n)\) (INT, V, p. 71, § 6, no. 2, prop. 3). The map \(h_n\) from X into \(\mathbf{C}\) such that \(h_n(x) = 0\) if \(x\notin \mathrm{U}-\mathrm{V}_n\) and \(h_n(x) = f(g(x))\) if \(x\in \mathrm{U}-\mathrm{V}_n\) is therefore \(\mu\)-measurable (INT, IV, p. 193, § 5, no. 10, prop. 16).

Since \(h_{n}(x) \to h(x)\) for \(\mu\)-almost every \(x \in U\), the restriction of \(h\) to U is \(\mu\)-measurable (INT, IV, p. 175, § 5, no. 4, th. 2). One concludes that \(h\) is \(\mu\)-measurable by the localization principle (INT, IV, p. 171, § 4, no. 2, prop. 4).

Remark. --- Under the hypotheses of the lemma, let us denote, by abuse of language, \(f \circ g\) the function \(h\) defined in the lemma.

If \(g = g_{1}\) locally \(\mu\)-almost everywhere on X, the functions \(g\) and \(g_{1}\) have the same \(\mu\)-essential image and one has \(f \circ g = f \circ g_{1}\) .

If \(g\) is a \(\mu\)-measurable function defined \(\mu\)-almost everywhere on X, we shall also denote by \(f \circ g\) the function \(f \circ g_{1}\), where \(g_{1}\) is a \(\mu\)-measurable function defined on X equal to g locally \(\mu\)-almost everywhere.

